\subsection{Assignment 17 -- Variazione percentuale annuale degli stream per categoria}

L’Assignment 17 richiedeva di analizzare la variazione percentuale, in aumento o
in diminuzione, del numero di stream per ciascuna categoria musicale rispetto
all’anno precedente.

La query MDX è stata progettata per confrontare il volume di stream di ogni
categoria musicale (\texttt{[Song].[Category]}) nell’anno corrente con quello
registrato nello stesso periodo dell’anno precedente. A tal fine, sono stati
utilizzati i seguenti costrutti.

La funzione \textit{PARALLELPERIOD} è stata impiegata per navigare la gerarchia
temporale e recuperare i valori relativi all’anno precedente, mantenendo
invariata la granularità dell’analisi. È stata inoltre definita una misura
calcolata, denominata \textit{Percentage Variation}, che automatizza il calcolo
della variazione percentuale degli stream, includendo una gestione esplicita dei
casi di divisione per zero tramite la funzione \textit{IIF}.

Nella formulazione iniziale, la query produceva un numero elevato di valori
\textit{NULL} o infiniti, rendendo i risultati difficili da interpretare. Per
migliorare la qualità dell’output, la soluzione è stata ottimizzata applicando
una clausola \textit{FILTER} combinata con \textit{NOT IsEmpty}, così da rimuovere
tutte le righe per cui non era possibile effettuare un confronto significativo.
In particolare, sono stati esclusi i casi in cui non esiste un anno precedente
di riferimento, come il primo anno presente nel database.

Questo filtraggio consente di isolare esclusivamente i record per i quali il
confronto matematico è effettivamente definito. Di conseguenza, una variazione
pari a $-100\%$ indica che una categoria musicale era presente nell’anno
precedente ma non registra stream nell’anno corrente, rappresentando
un’informazione di business reale. In assenza di tale filtraggio, non sarebbe
stato possibile distinguere tra una categoria effettivamente scomparsa e una per
cui mancano semplicemente dati storici.

Infine, l’ordinamento basato sulla \textit{Member\_Key} della dimensione temporale
assicura una progressione cronologica coerente, facilitando la lettura e
l’interpretazione dei trend storici.

Dall’output ottenuto emergono indicazioni rilevanti dal punto di vista del
business. Variazioni percentuali fortemente positive segnalano l’emergere o
l’esplosione di nuovi trend di mercato, mentre variazioni negative evidenziano
fasi di contrazione che possono richiedere analisi strategiche dedicate. Un calo
generalizzato nell’ultimo anno disponibile costituisce infine un indicatore
tecnico della possibile incompletezza dei dati, dimostrando come la query sia
sensibile alla completezza temporale del dataset.
